% Arbitrering mellan två noder: A förlorar vid bit 8.
\begin{tikzpicture}[
  every node/.style={font=\scriptsize},
  bit/.style={minimum width=6.2mm, minimum height=5mm, inner sep=0pt, align=center,
              font=\scriptsize}]

  \def\bits{{0,1,0,1,0,1,0,0,1,1,1}}   % A: 0x2A7
  \def\bitsb{{0,1,0,1,0,1,0,0,0,1,1}}  % B: 0x2A3

  % Rubriker.
  \node[anchor=east, font=\scriptsize\bfseries] at (-0.15,0.6)  {bit};
  \node[anchor=east, font=\scriptsize\bfseries] at (-0.15,0)    {A};
  \node[anchor=east, font=\scriptsize\bfseries] at (-0.15,-0.6) {B};
  \node[anchor=east, font=\scriptsize\bfseries] at (-0.15,-1.2) {buss};

  \foreach \i in {0,...,10} {
    \pgfmathsetmacro{\x}{\i*0.65}
    \pgfmathtruncatemacro{\n}{\i+1}
    \pgfmathtruncatemacro{\a}{\bits[\i]}
    \pgfmathtruncatemacro{\b}{\bitsb[\i]}
    \pgfmathtruncatemacro{\bus}{min(\a,\b)}

    \node[bit, muted] at (\x,0.6) {\n};

    % A slutar sända efter bit 8 (index 7).
    \ifnum\i<9
      \node[bit, draw=rule, fill=white] at (\x,0) {\a};
    \else
      \node[bit, draw=rule, fill=fillmuted] at (\x,0) {\textcolor{ink}{\textbf{--}}};
    \fi

    \node[bit, draw=rule, fill=white] at (\x,-0.6) {\b};

    \ifnum\bus=0
      \node[bit, draw=rule, fill=fillaccent!70] at (\x,-1.2) {\bus};
    \else
      \node[bit, draw=rule, fill=fillaccent2!60] at (\x,-1.2) {\bus};
    \fi
  }

  % Markera bit 8.
  \pgfmathsetmacro{\xl}{8*0.65}
  \draw[accent, line width=0.8pt] (\xl-0.33,0.32) rectangle (\xl+0.33,-1.45);
  \node[accent, anchor=north, font=\scriptsize\bfseries, align=center]
    at (\xl,-1.62) {A sänder 1, läser 0\\[-0.2ex] A förlorar};

\end{tikzpicture}
